\section*{Bab \ref{ch:b1:gene-expression} --- Pengungkapan Gen: Transkripsi dan Translasi}

\begin{solution}{exo:b1:gene-expression:1}
\omterm{def:b1:nucleic-acids:chain}{RNA} $5'$-AUGCCUAAG-$3'$; untai penyandinya $5'$-ATGCCTAAG-$3'$.
\end{solution}

\begin{solution}{exo:b1:gene-expression:2}
Sebuah tudung $5'$ (guanina yang diubah), sebuah \omterm{prop:b1:gene-expression:processing}{ekor poli-A} $3'$, dan
penyingkiran \omterm{def:b1:genomes:gene}{intronnya} lewat penyambungan.
\end{solution}

\begin{solution}{exo:b1:gene-expression:3}
AUG CCG AAA GUU UGA: Met-Pro-Lys-Val, lalu berhenti.
\end{solution}

\begin{solution}{exo:b1:gene-expression:4}
A: tRNA bermuatan yang masuk diizinkan lalu diperiksa. P: tRNA yang
membawa rantai yang tumbuh; ikatannya terbentuk antara rantai itu dan
\omterm{def:b1:proteins:aminoacid}{asam amino} sisi A-nya. E: tRNA yang telah dikosongkan dalam perjalanan
keluarnya.
\end{solution}

\begin{solution}{exo:b1:gene-expression:5}
Yang menyandi $2400 - 600 = 1800$ \omterm{def:b1:nucleic-acids:nucleotide}{nukleotida}: 599 residu (600 \omterm{def:b1:gene-expression:code}{kodon},
satu di antaranya berhenti). Waktunya $599/4 = \qty{150}{s}$. \omterm{def:b1:gene-expression:ribosome}{Ribosomnya}
$1800/90 = 20$ sekaligus.
\end{solution}

\begin{solution}{exo:b1:gene-expression:6}
Pesannya dibaca bertiga dari sebuah awal yang tetap; satu atau dua basa
tambahan menggeser tiap \omterm{def:b1:gene-expression:code}{kodon} di hilirnya dan sisa \omterm{def:b1:proteins:peptide}{proteinnya} menjadi
omong kosong, yang biasanya berakhir pada sebuah berhenti yang terlalu
dini. Tiga basa tambahan menambahkan satu \omterm{def:b1:gene-expression:code}{kodon} lalu memulihkan
kerangkanya: \omterm{def:b1:proteins:peptide}{proteinnya} memiliki satu residu tambahan dan beberapa
residu yang keliru di antara penyisipannya, dan sering masih bekerja.
\end{solution}

\begin{solution}{exo:b1:gene-expression:7}
CAG (Gln) menjadi UAG (berhenti): rantainya berakhir pada residu 150,
sebuah \omterm{def:b1:proteins:peptide}{protein} yang terpangkas dan hampir pasti tak aktif (mutasi
nonsens). CAG menjadi CAA: tetap Gln, bisu. CAG menjadi CGG: Arg
menggantikan Gln, sebuah penggantian (missens) yang pengaruhnya
bergantung pada kedudukannya.
\end{solution}

\begin{solution}{exo:b1:gene-expression:8}
\omterm{def:b1:gene-expression:ribosome}{Ribosomnya} hanya memeriksa pemasangan \omterm{def:b1:gene-expression:code}{kodon} dengan \omterm{def:b1:gene-expression:trna}{antikodonnya}; ia
tak dapat melihat \omterm{def:b1:proteins:aminoacid}{asam aminonya}. Bila sisteina pada tRNA-nya diubah
secara kimia menjadi alanina, \omterm{def:b1:gene-expression:ribosome}{ribosomnya} menyisipkan alanina di mana
pun pesannya menyebut sisteina: tRNA-nya, bukan \omterm{def:b1:proteins:aminoacid}{asam aminonya}, yang
dibaca. Karena itu sintetasenya, yang memasangkan tiap \omterm{def:b1:proteins:aminoacid}{asam amino}
dengan tRNA yang tepat, adalah penerjemah yang sesungguhnya.
\end{solution}

\begin{solution}{exo:b1:gene-expression:9}
$400\times 4 = 1600$ ikatan; dua GTP \omterm{def:b1:gene-expression:ribosome}{ribosomnya} tiap residu adalah
separuhnya, dua milik sintetasenya separuh yang lain.
\end{solution}

\begin{solution}{exo:b1:gene-expression:10}
Pada eukariot transkripnya dibuat dalam \omterm{def:b1:cell-unit-of-life:prokeuk}{intinya} dan \omterm{def:b1:gene-expression:ribosome}{ribosomnya} berada
dalam \omterm{def:b1:eukaryotic-cell:organelle}{sitosolnya}, yang dipisahkan selubungnya; dan transkripnya belum
menjadi sebuah pembawa pesan sampai ia disambung dan ditudungi.
Pemisahan itu memungkinkan pemrosesannya sendiri --- penyambungan
berselang-selang, kendali ekspornya, dan sebuah pemeriksaan bahwa
pesannya lengkap sebelum ia dibaca.
\end{solution}

\begin{solution}{exo:b1:gene-expression:11}
Empat pembawa pesan: dengan kedua \omterm{def:b1:genomes:gene}{eksonnya} (190 \omterm{def:b1:nucleic-acids:nucleotide}{nukleotida}
ditambahkan), dengan 3 saja (90), dengan 4 saja (100), dengan tak satu
pun (0). Sebuah kerangka terjaga bila panjang yang ditambahkan
kelipatan 3: 0 dan 90 menjaganya; 100 dan 190 menggesernya, sehingga
pembawa pesan dengan \omterm{def:b1:genomes:gene}{ekson} 4 saja atau dengan kedua \omterm{def:b1:genomes:gene}{eksonnya} membaca
\omterm{def:b1:genomes:gene}{ekson} 5 di luar kerangkanya lalu memangkas \omterm{def:b1:proteins:peptide}{proteinnya}. Penyambungan
berselang-selang harus menghormati kerangkanya, dan karena itulah \omterm{def:b1:genomes:gene}{ekson}
sering kelipatan tiga.
\end{solution}

\begin{solution}{exo:b1:gene-expression:12}
Sewenang-wenang: tak ada dalam kimianya yang mengatakan bahwa UUU harus
berarti Phe; penetapannya adalah kesepakatan yang dipegang
sintetasenya. Dioptimumkan: kedudukan ketiganya paling merosot,
sehingga kekeliruan dan mutasi yang jatuh di sana sering bisu, dan
\omterm{def:b1:gene-expression:code}{kodon} yang berbeda satu basa cenderung menyandi \omterm{def:b1:proteins:aminoacid}{asam amino} yang serupa,
sehingga sebuah penggantian sering ringan --- kodenya memperkecil
kerusakan akibat kekeliruan. Membeku: sebuah sintetase yang mengubah
kekhasannya akan mengubah tiap \omterm{def:b1:proteins:peptide}{protein} \omterm{def:b1:cell-unit-of-life:cell}{selnya} sekaligus, ribuan
jumlahnya, pada saat yang sama; nyaris tak ada \omterm{def:b1:cell-unit-of-life:cell}{sel} yang selamat dari
itu, sehingga kodenya tak dapat melayang dan telah tetap sejak leluhur
bersama semua makhluk hidup.
\end{solution}

\begin{solution}{pb:b1:gene-expression:1}
\textbf{1.} $30\,000/30 = \qty{1000}{s}$, sekitar \qty{17}{min}.
\textbf{2.} $28\,000/30\,000 = \qty{93}{\%}$.
\textbf{3.} $30\,000/2000 = 15$.
\textbf{4.} $2\times 30\,000 = \num{60000}$ \omterm{def:b1:water-small-molecules:atp}{ATP}.
\textbf{5.} $1000/10 = 100$ polimerase pada \omterm{def:b1:genomes:gene}{gennya}; 360 pembawa pesan
tiap jam.
\textbf{6.} Tujuh \omterm{def:b1:genomes:gene}{intron}, rata-rata $28\,000/7 = \qty{4000}{}$ \omterm{def:b1:nucleic-acids:nucleotide}{nukleotida}.
\textbf{7.} Transkripsinya (\qty{17}{min}); \omterm{def:b1:genomes:gene}{intronnya} disambung seiring
dibuatnya, dan yang terakhir hanya menambahkan semenit.
\textbf{8.} \omterm{def:b1:genomes:gene}{Gennya} $30\,000\times 0.34\,\mathrm{nm} = \qty{10}{\micro m}$; pembawa pesannya \qty{0.68}{\micro m}.
\textbf{9.} $500/4 = \qty{125}{s}$.
\textbf{10.} $2000/80 = 25$.
\textbf{11.} Satu rantai selesai tiap $80/4 = \qty{20}{s}$: 180 tiap
jam.
\textbf{12.} $2/\ln 2 = \qty{2.9}{h}$ produksi penuh: sekitar 520
\omterm{def:b1:proteins:peptide}{protein}.
\textbf{13.} $4\times 500 = 2000$ \omterm{def:b1:water-small-molecules:atp}{ATP}.
\textbf{14.} $60\,000/520 = 115$ \omterm{def:b1:water-small-molecules:atp}{ATP} tiap \omterm{def:b1:proteins:peptide}{protein}.
\textbf{15.} Sekitar 2100 \omterm{def:b1:water-small-molecules:atp}{ATP}, \qty{95}{\%} di antaranya dalam translasi.
\textbf{16.} $1 - (1 - 10^{-5})^{1503} \approx 1503\times 10^{-5} =
\qty{1.5}{\%}$.
\textbf{17.} $1 - (1 - 10^{-4})^{500} \approx \qty{5}{\%}$.
\textbf{18.} Transkripsi: $1.5\%\times 0.75\times 0.67 \approx
0.75\%$ pembawa pesannya cacat, dan dengan itu \omterm{def:b1:proteins:peptide}{proteinnya}; translasi:
$5\%\times 0.67 = 3.3\%$; sekitar \qty{4}{\%} molekulnya.
\textbf{19.} Sebuah kekeliruan \omterm{def:b1:proteins:peptide}{protein} terkurung pada satu molekul,
yang segera diuraikan; sebuah kekeliruan pembawa pesan disalin ke dalam
ratusan \omterm{def:b1:proteins:peptide}{protein}; sebuah kekeliruan replikasi diwarisi tiap
keturunannya selamanya. Biaya sebuah kekeliruan, dan dengan itu
ketepatan yang layak dibayar, naik pada tiap langkah mundur.
\textbf{20.} $3000/180 = 17$ pembawa pesan yang hadir.
\textbf{21.} $17\times 25 = 420$ \omterm{def:b1:gene-expression:ribosome}{ribosom}.
\textbf{22.} $\num{3e9}\times 4 = \num{1.2e10}$ \omterm{def:b1:water-small-molecules:atp}{ATP} tiap hari
$= \qty{2e-14}{mol}$, \qty{1e-9}{J} tiap hari, yakni \qty{1.2e-14}{W}
--- \qty{6}{fW} tiap mikrometer kubik.
\textbf{23.} $10^9$ \omterm{def:b1:water-small-molecules:atp}{ATP} tiap sekon adalah $\num{8.6e13}$ tiap hari:
sintesis \omterm{def:b1:proteins:peptide}{proteinnya} seperseratusnya pada \omterm{def:b1:cell-unit-of-life:cell}{sel} yang memperbarui diri
perlahan ini (pada sebuah bakteri yang membelah ia separuhnya).
\textbf{24.} Tak ada pembawa pesan yang matang: transkrip primernya
menumpuk dalam \omterm{def:b1:cell-unit-of-life:prokeuk}{intinya} dan \omterm{def:b1:proteins:peptide}{proteinnya} lenyap seiring meluruhnya
pembawa pesannya, dalam beberapa jam. \omterm{def:b1:proteins:peptide}{Protein} bakterinya, yang \omterm{def:b1:genomes:gene}{gennya}
tak berintron, tak terpengaruh.
\textbf{25.} Sekitar 2100 \omterm{def:b1:water-small-molecules:atp}{ATP} tiap molekul --- 2000 bagi translasinya,
seratus bagi bagiannya dari transkripnya; \omterm{def:b1:proteins:peptide}{protein} pertama sesudah
\qty{1000}{s} transkripsi, semenit pemrosesan dan ekspor, serta
\qty{125}{s} translasi: sekitar dua puluh menit.
\end{solution}
